\section{Dlaczego nie przeczy to objętości?}
\subsection{Miara jest określona tylko na pewnych zbiorach}
Miara Lebesgue'a w $\R^3$ przypisuje zbiorom mierzalnym objętość $\lambda(E)\in[0,\infty]$. Do poniższego argumentu potrzebne są tylko jej standardowe własności:
\begin{itemize}
\item kule są mierzalne i $0<\lambda(B^3)<\infty$;
\item izometrie zachowują mierzalność i objętość;
\item skończona unia rozłącznych zbiorów mierzalnych ma objętość równą sumie ich objętości.
\end{itemize}
W standardowej normalizacji $\lambda(B^3)=4\pi/3$, ale konkretna liczba nie jest istotna. Twierdzenie o istnieniu i własnościach miary Lebesgue'a jest tu wyłącznie tłem interpretacyjnym; główny dowód w poprzednich sekcjach nie korzystał z teorii miary.
\begin{stwierdzenie}[Co najmniej jeden kawałek jest niemierzalny]\label{st:niemierzalnosc}
W rozkładzie z twierdzenia~\ref{tw:banach-tarski} co najmniej jeden kawałek jest niemierzalny w sensie Lebesgue'a.
\end{stwierdzenie}
\begin{proof}
Załóżmy przeciwnie, że wszystkie $P_i$ są mierzalne. Ponieważ tworzą partycję $B^3$, mamy $\lambda(B^3)=\sum_{i=1}^{N}\lambda(P_i)$. Obrazy $f_i(P_i)$ też są mierzalne, mają te same objętości i tworzą partycję dwóch rozłącznych kul przystających do $B^3$. Zatem
\[
 2\lambda(B^3)=\sum_{i=1}^{N}\lambda(f_i(P_i))=
 \sum_{i=1}^{N}\lambda(P_i)=\lambda(B^3).
\]
Jest to niemożliwe dla dodatniej skończonej liczby $\lambda(B^3)$. Sprzeczne było założenie o mierzalności wszystkich części, a nie twierdzenie Banacha--Tarskiego.
\end{proof}
\begin{uwaga}
Sformułowanie „objętość wzrosła dwukrotnie podczas obracania kawałków” jest mylące. Co najmniej jednemu kawałkowi nie można przypisać objętości Lebesgue'a, więc nie można zsumować objętości wszystkich części. Nie mamy prawa zastosować addytywności miary do dowolnych podzbiorów tylko dlatego, że ich unia jest kulą.
\end{uwaga}

\subsection{Co jest, a co nie jest konstruktywne?}
Generatory obrotów są podane jawnymi macierzami. Cztery zbiory słów są zdefiniowane jawnie. Wzory na składanie kawałków też są jawne. Jednak zbiór reprezentantów $T$ został wybrany aksjomatem wyboru. Nie otrzymaliśmy algorytmu, który dla dowolnego punktu kuli rozstrzyga, do którego kawałka należy.

Nie twierdzimy też, że twierdzenie Banacha--Tarskiego jest równoważne pełnemu aksjomatowi wyboru ani że każdy fragment użytego wyboru jest logicznie niezbędny. Wykazaliśmy konkretnie: w ZFC istnieją zbiory i ruchy o opisanych własnościach. Taka jest treść tego dowodu.

Kula jako model ciągłego zbioru punktów nie jest fizycznym przedmiotem zbudowanym z atomów. Rozkład nie opisuje procedury możliwej do wykonania nożem ani technologii powielania materii. Kawałki nie muszą mieć wnętrza ani regularnego brzegu.

